\newpage
\section*{Solucionario completo}

\noindent Cada respuesta incluye la clave oficial y su desarrollo económico y analítico detallado.

\subsection*{Bloque A. Fundamentos teóricos y conceptuales}

\sol{1}{B}{\textbf{Tema:} Estrategias dominantes. \textbf{Explicación:} Una estrategia es estrictamente dominante para el jugador $i$ si para cualquier perfil de estrategias de los rivales $s_{-i}$, se cumple $u_i(s_i, s_{-i}) > u_i(s_i', s_{-i})$ para todo $s_i' 
eq s_i$. No depende de lo que elijan los demás.}

\sol{2}{A}{\textbf{Tema:} Equilibrio de Nash. \textbf{Explicación:} En el equilibrio de Nash, la estrategia de cada jugador es una mejor respuesta a las estrategias de equilibrio de los demás ($s_i^* \in BR_i(s_{-i}^*)$). Ningún jugador puede mejorar su pago cambiando unilateralmente de acción.}

\sol{3}{D}{\textbf{Tema:} Dilema del prisionero. \textbf{Explicación:} En el dilema del prisionero, confesar (desviarse) es la estrategia estrictamente dominante para ambos jugadores. El equilibrio de Nash resultante (Confesar, Confesar) genera pagos estrictamente inferiores para ambos que el resultado cooperativo (No confesar, No confesar), por lo que es ineficiente en el sentido de Pareto.}

\sol{4}{B}{\textbf{Tema:} Juegos dinámicos e información. \textbf{Explicación:} En información perfecta, todo conjunto de información contiene exactamente un único nodo del árbol de juego, lo que significa que en cada momento el jugador de turno conoce con total certeza la historia completa del juego previa a su decisión.}

\sol{5}{A}{\textbf{Tema:} Inducción hacia atrás. \textbf{Explicación:} La inducción hacia atrás comienza identificando las acciones óptimas en los nodos de decisión finales (aquellos inmediatamente anteriores a los pagos terminales) y retrocede secuencialmente hacia la raíz del árbol, garantizando la racionalidad secuencial.}

\sol{6}{B}{\textbf{Tema:} Equilibrio perfecto en subjuegos. \textbf{Explicación:} El SPNE exige que las estrategias constituyan un equilibrio de Nash en todos y cada uno de los subjuegos del juego original. Esto elimina los equilibrios sustentados en amenazas no creíbles que un jugador no ejecutaría si llegara al nodo correspondiente.}

\sol{7}{A}{\textbf{Tema:} Estrategias mixtas. \textbf{Explicación:} Para que un jugador esté dispuesto a aleatorizar entre varias estrategias puras, todas ellas deben generarle exactamente el mismo pago esperado en equilibrio, y dicho pago esperado debe ser mayor o igual al de cualquier estrategia pura con probabilidad cero.}

\sol{8}{B}{\textbf{Tema:} Teorema Popular (Folk Theorem). \textbf{Explicación:} El Folk Theorem demuestra que en juegos repetidos indefinidamente con un factor de descuento intertemporal $\delta$ cercano a 1, la amenaza de castigos creíbles futuros compensa las ganancias de corto plazo de la desviación, permitiendo sostener infinitos equilibrios cooperativos.}

\sol{9}{B}{\textbf{Tema:} Modelo de Cournot. \textbf{Explicación:} Cournot modela la competencia simultánea en cantidades con productos homogéneos. El precio de mercado resultante se ajusta a través de la función inversa de demanda para vaciar el mercado ante la producción agregada elegida por las firmas.}

\sol{10}{D}{\textbf{Tema:} Funciones de reacción en Cournot. \textbf{Explicación:} La función de reacción $q_i = BR_i(q_{-i})$ se obtiene de maximizar el beneficio de la firma $i$ respecto a $q_i$. Con demanda lineal y costos constantes, tiene pendiente negativa, indicando que a mayor producción del rival, menor es la cantidad óptima que la firma debe producir.}

\sol{11}{C}{\textbf{Tema:} Modelo de Bertrand homogéneo. \textbf{Explicación:} En Bertrand con producto homogéneo e información perfecta, los consumidores compran exclusivamente a la empresa con el menor precio. Si fijan el mismo precio, se asume usualmente que se reparten la demanda en partes iguales.}

\sol{12}{C}{\textbf{Tema:} Paradoja de Bertrand. \textbf{Explicación:} Bastan solo dos empresas compitiendo en precios con costos marginales constantes, producto idéntico y capacidad ilimitada para que la guerra de precios reduzca el precio hasta $P = CMg$, disipando por completo las rentas oligopólicas como si hubiera competencia perfecta.}

\sol{13}{B}{\textbf{Tema:} Modelo de Bertrand-Edgeworth. \textbf{Explicación:} Edgeworth demostró que si las empresas enfrentan restricciones de capacidad que les impiden abastecer toda la demanda a $P = CMg$, ninguna empresa tiene incentivos para bajar el precio hasta el costo marginal, rompiéndose la paradoja y generando ciclos o equilibrios en estrategias mixtas.}

\sol{14}{D}{\textbf{Tema:} Bertrand con diferenciación. \textbf{Explicación:} Con bienes diferenciados, si el rival sube su precio, parte de su clientela se traslada a la empresa, permitiéndole a esta subir también su propio precio para maximizar beneficios. Por tanto, $\partial p_i / \partial p_j > 0$ (complementos estratégicos).}

\sol{15}{A}{\textbf{Tema:} Sustitutos estratégicos. \textbf{Explicación:} Las acciones son sustitutos estratégicos si un aumento en la acción de la firma 1 reduce el beneficio marginal de la firma 2, haciendo que esta reaccione reduciendo su acción (pendiente negativa de la función de mejor respuesta, como en Cournot).}

\sol{16}{A}{\textbf{Tema:} Modelo de Stackelberg. \textbf{Explicación:} En Stackelberg, el líder posee ventaja temporal: maximiza su función de beneficios sustituyendo formalmente la función de reacción de la seguidora $q_S(q_L)$, eligiendo una producción mayor que en Cournot.}

\sol{17}{D}{\textbf{Tema:} Compromiso creíble en Stackelberg. \textbf{Explicación:} Si la decisión de producción del líder fuera fácilmente revocable o modificable sin costo, la seguidora anticiparía que el líder reajustaría su cantidad, colapsando el juego nuevamente al equilibrio simultáneo de Cournot.}

\sol{18}{C}{\textbf{Tema:} Comparación de cantidades en oligopolio. \textbf{Explicación:} Para demanda lineal: $Q_M = \frac{1}{2}Q_{CP}$, en duopolio de Cournot $Q_C = \frac{2}{3}Q_{CP}$, en Stackelberg $Q_S = \frac{3}{4}Q_{CP}$, y en competencia perfecta (o Bertrand homogéneo) $Q_{CP} = \frac{a-c}{b}$.}

\sol{19}{B}{\textbf{Tema:} Comparación de precios en oligopolio. \textbf{Explicación:} Como el precio es función decreciente de la cantidad agregada, a mayor cantidad producida menor precio. Por tanto, el precio más bajo es el de competencia perfecta ($P=CMg$), seguido de Stackelberg, luego Cournot y el más alto es el monopolio.}

\sol{20}{C}{\textbf{Tema:} Incentivo a la desviación en cárteles. \textbf{Explicación:} En la cuota colusoria, el precio de mercado es superior al costo marginal ($P > CMg$). Una empresa individual puede obtener un beneficio extraordinario mayor desviándose unilateralmente (haciendo trampa) y aumentando su producción mientras las demás restringen la suya.}

\sol{21}{C}{\textbf{Tema:} Estrategia del gatillo (Grim Trigger). \textbf{Explicación:} La condición de no desvío exige: $\frac{\pi_{cartel}}{1-\delta} \ge \pi_{desvio} + \frac{\delta}{1-\delta} \pi_{castigo}$. Si el factor de descuento $\delta$ es lo bastante alto, la pérdida por recibir para siempre el beneficio competitivo o de Cournot supera la ganancia transitoria del desvío.}

\sol{22}{A}{\textbf{Tema:} Número de firmas y colusión. \textbf{Explicación:} A mayor número de empresas, la ganancia que le corresponde a cada una dentro del cártel ($\Pi_M/n$) es menor, mientras que el incentivo relativo a hacer trampa y capturar una tajada mayor se mantiene alto. Por tanto, se requiere un $\delta^*$ mayor para sostener el cártel.}

\sol{23}{A}{\textbf{Tema:} Competencia Monopolística. \textbf{Explicación:} La competencia monopolística combina elementos de monopolio (poder de mercado sobre una variedad diferenciada con pendiente negativa de demanda) y de competencia perfecta (muchos participantes y libre entrada a largo plazo).}

\sol{24}{B}{\textbf{Tema:} Modelo de Hotelling. \textbf{Explicación:} En el modelo de Hotelling, las diferencias entre bienes se basan en atributos o localización espacial donde no existe un producto objetivamente mejor para todos los consumidores; la elección depende del costo de transporte o preferencia subjetiva (diferenciación horizontal).}

\sol{25}{B}{\textbf{Tema:} Diferenciación vertical. \textbf{Explicación:} La diferenciación vertical implica que existe un ordenamiento unánime de calidad: si dos variedades tuvieran el mismo precio, ningún consumidor compraría la de menor calidad. Las elecciones difieren solo porque los consumidores difieren en su ingreso o disposición a pagar por calidad.}

\sol{26}{A}{\textbf{Tema:} Equilibrio de largo plazo de Chamberlin. \textbf{Explicación:} La libre entrada de empresas con variedades sustitutas desplaza la demanda de cada firma hacia adentro hasta que los beneficios económicos se anulan ($P = CMe$). La tangencia se da en el tramo decreciente del $CMe$ porque la curva de demanda tiene pendiente negativa.}

\sol{27}{A}{\textbf{Tema:} Exceso de capacidad. \textbf{Explicación:} Dado que la condición de tangencia entre la demanda con pendiente negativa y la curva de $CMe$ ocurre en el tramo decreciente de esta última, la empresa produce menos de lo que exigiría el punto de costo medio mínimo (escala eficiente), existiendo exceso de capacidad.}

\sol{28}{D}{\textbf{Tema:} Bienestar y variedad de productos. \textbf{Explicación:} Aunque en competencia monopolística $P > CMg$ (lo que genera pérdida irrecuperable de eficiencia en costos), los consumidores se benefician de disponer de múltiples variedades diferenciadas que se ajustan mejor a sus preferencias heterogéneas.}

\sol{29}{A}{\textbf{Tema:} Publicidad informativa vs persuasiva. \textbf{Explicación:} La publicidad informativa aporta datos transparentes sobre disponibilidad, características y precios, incrementando la rivalidad y la elasticidad de la demanda; la persuasiva busca fidelizar clientes, aumentando la diferenciación percibida y el margen de mercado.}

\sol{30}{A}{\textbf{Tema:} Estabilidad en Cournot. \textbf{Explicación:} La condición de estabilidad estandar exige $|\partial BR_1 / \partial q_2| \cdot |\partial BR_2 / \partial q_1| < 1$, lo que en el caso simétrico equivale a que la pendiente de la función de mejor respuesta en el espacio $(q_1, q_2)$ tenga un valor absoluto menor a la unidad.}

\subsection*{Bloque B. Aplicaciones cuantitativas y optimización}

\sol{31}{D}{\textbf{Tema:} Cálculo de Cournot simétrico. \textbf{Explicación:} La firma 1 maximiza $\pi_1 = (100 - q_1 - q_2 - 10)q_1$. Condición de primer orden: $90 - 2q_1 - q_2 = 0 \implies q_1 = (90 - q_2)/2$. Por simetría $q_1 = q_2 = q \implies 3q = 90 \implies q^* = 30$.}

\sol{32}{A}{\textbf{Tema:} Precio y beneficios en Cournot. \textbf{Explicación:} Producción agregada: $Q = q_1 + q_2 = 30 + 30 = 60$. Precio: $P = 100 - 60 = 40$. Beneficio de cada firma: $\pi_i = (P - c)q_i = (40 - 10) \times 30 = 30 \times 30 = 900$.}

\sol{33}{C}{\textbf{Tema:} Cournot con n empresas. \textbf{Explicación:} Cada empresa produce $q_i = \frac{a-c}{b(n+1)}$. Sumando las $n$ firmas: $Q^* = \sum q_i = \frac{n(a-c)}{b(n+1)}$.}

\sol{34}{B}{\textbf{Tema:} Precio en oligopolio de n firmas. \textbf{Explicación:} $q_i = (150 - 30)/(1 \times (5+1)) = 120 / 6 = 20$. Producción total: $Q = 5 \times 20 = 100$. Precio: $P = 150 - 100 = 50$.}

\sol{35}{D}{\textbf{Tema:} Cournot con costos asimétricos. \textbf{Explicación:} $BR_1: q_1 = (100 - q_2)/2$; $BR_2: q_2 = (80 - q_1)/2$. Sustituyendo: $2(80 - q_1) = 80 - q_1$, es decir $q_1 = 50 - 0.5(40 - 0.5 q_1) = 30 + 0.25 q_1 \implies 0.75 q_1 = 30 \implies q_1^* = 40$. Luego $q_2^* = (80 - 40)/2 = 20$.}

\sol{36}{A}{\textbf{Tema:} Beneficios con costos asimétricos. \textbf{Explicación:} $\pi_1^* = (P - c_1)q_1^* = (60 - 20) \times 40 = 40 \times 40 = 1600$. $\pi_2^* = (P - c_2)q_2^* = (60 - 40) \times 20 = 20 \times 20 = 400$. La empresa más eficiente cuadruplica el beneficio de la menos eficiente.}

\sol{37}{C}{\textbf{Tema:} Líder de Stackelberg. \textbf{Explicación:} La seguidora reacciona con $q_2 = (90 - q_1)/2$. El líder sustituye en su beneficio: $\pi_1 = [100 - q_1 - (90-q_1)/2 - 10]q_1 = (45 - 0.5 q_1)q_1$. Derivando: $45 - q_1 = 0 \implies q_L^* = 45$.}

\sol{38}{D}{\textbf{Tema:} Seguidora y precio en Stackelberg. \textbf{Explicación:} $q_S^* = (90 - 45)/2 = 22.5$. Producción total: $Q = 45 + 22.5 = 67.5$. Precio: $P = 100 - 67.5 = 32.5$.}

\sol{39}{C}{\textbf{Tema:} Beneficios en Stackelberg. \textbf{Explicación:} $\pi_L = (32.5 - 10) \times 45 = 22.5 \times 45 = 1012.5$. $\pi_S = (32.5 - 10) \times 22.5 = 22.5 \times 22.5 = 506.25$. El líder gana más que en Cournot (1012.5 vs 900) y la seguidora menos (506.25 vs 900).}

\sol{40}{B}{\textbf{Tema:} Equilibrio de Bertrand homogéneo. \textbf{Explicación:} Por la paradoja de Bertrand, la competencia en precios induce a subcotizaciones sucesivas hasta que $P = c_1 = c_2 = 30$. A este precio, $Q = 150$, y los beneficios económicos extraordinarios son nulos.}

\sol{41}{B}{\textbf{Tema:} Bertrand con costos asimétricos. \textbf{Explicación:} La empresa 1 (más eficiente) puede fijar un precio justo por debajo del costo marginal de la empresa 2 ($P = 35 - \epsilon \approx 34.99$). A ese precio captura toda la demanda: $Q = 100 - 35 = 65$, obteniendo un margen de $35 - 25 = 10$ por unidad y un beneficio $\pi_1 \approx 10 \times 65 = 650$.}

\sol{42}{D}{\textbf{Tema:} Reacción en Bertrand diferenciado. \textbf{Explicación:} $\pi_1 = p_1(24 - 2p_1 + p_2) = 24p_1 - 2p_1^2 + p_1 p_2$. Condición de primer orden: $\frac{\partial \pi_1}{\partial p_1} = 24 - 4p_1 + p_2 = 0 \implies 4p_1 = 24 + p_2 \implies p_1 = 6 + 0.25 p_2$.}

\sol{43}{D}{\textbf{Tema:} Equilibrio simétrico en Bertrand diferenciado. \textbf{Explicación:} Por simetría: $p = 6 + 0.25 p \implies 0.75 p = 6 \implies p^* = 8$. Cantidad demandada: $q_i = 24 - 2(8) + 8 = 16$. Beneficio: $\pi_i = 8 \times 16 = 128$.}

\sol{44}{B}{\textbf{Tema:} Cártel colusorio simétrico. \textbf{Explicación:} El cártel actúa como monopolio: $\max (140 - Q - 20)Q = (120 - Q)Q \implies IMg = 120 - 2Q = 0 \implies Q_M = 60$. Precio: $P = 140 - 60 = 80$. Beneficio total: $\Pi_M = (80-20) \times 60 = 3600$. Cada empresa produce $q_i = 30$ y gana $\pi_i^M = 1800$.}

\sol{45}{B}{\textbf{Tema:} Desviación unilateral del cártel. \textbf{Explicación:} La empresa 1 optimiza según su función de reacción de Cournot dada $q_2 = 30$: $q_1 = (120 - 30)/2 = 45$. Producción agregada con trampa: $Q = 45 + 30 = 75$. Precio: $P = 140 - 75 = 65$. Beneficio con desvío: $\pi_1^{dev} = (65 - 20) \times 45 = 45 \times 45 = 2025$.}

\sol{46}{C}{\textbf{Tema:} Factor de descuento crítico en Cournot. \textbf{Explicación:} La condición de sostenibilidad es $\delta \ge \frac{\pi_{dev} - \pi_{cartel}}{\pi_{dev} - \pi_C} = \frac{2025 - 1800}{2025 - 1600} = \frac{225}{425} = \frac{9}{17} \approx 0.5294$.}

\sol{47}{B}{\textbf{Tema:} Factor de descuento en Bertrand con n firmas. \textbf{Explicación:} En Bertrand con $n$ firmas: $\frac{\Pi_M/n}{1-\delta} \ge \Pi_M + 0 \implies \frac{1}{n(1-\delta)} \ge 1 \implies 1 - \delta \le \frac{1}{n} \implies \delta \ge 1 - \frac{1}{n}$. Para $n=4$: $\delta^* = 1 - 1/4 = 3/4 = 0.75$.}

\sol{48}{C}{\textbf{Tema:} Nash en dilema de precios 2x2. \textbf{Explicación:} Para la Firma 1: si F2 juega PA, $140 > 100$ (elige PB); si F2 juega PB, $50 > 20$ (elige PB). PB es estrictamente dominante para ambas. El único equilibrio de Nash es (PB, PB) con pagos (50, 50).}

\sol{49}{D}{\textbf{Tema:} Equilibrio mixto en coordinación. \textbf{Explicación:} Para que J2 sea indiferente, sus pagos esperados deben ser iguales: $E[u_2(A)] = E[u_2(B)] \implies p(1) + (1-p)(0) = p(0) + (1-p)(2) \implies p = 2 - 2p \implies 3p = 2 \implies p = 2/3$.}

\sol{50}{D}{\textbf{Tema:} SPNE de disuasión de entrada. \textbf{Explicación:} Por inducción hacia atrás: en el subjuego tras la entrada, el Incumbente compara $u_I(A) = 40$ con $u_I(L) = 20$. Como $40 > 20$, elige Acomodar. El Entrante anticipa esto y compara $u_E(E) = 15$ con $u_E(NE) = 0$. Como $15 > 0$, el Entrante entra. El SPNE es (E; Acomoda) con pagos (15, 40).}

\sol{51}{C}{\textbf{Tema:} Competencia monopolística a corto plazo. \textbf{Explicación:} $\pi_i = p_i(60 - 2p_i + 2\bar{p}) - 10(60 - 2p_i + 2\bar{p})$. $\frac{\partial \pi_i}{\partial p_i} = 60 - 4p_i + 2\bar{p} + 20 = 80 - 4p_i + 2\bar{p} = 0$. En simetría $p_i = \bar{p} = p^* \implies 80 - 4p^* + 2p^* = 0 \implies 2p^* = 80 \implies p^* = 40$... corrección: con $CT=10q$, $IMg = CMg$: $80-4p+2p=0 \implies p=40$ si fuera 80, pero aquí es: $q_i = 60 - 2p_i + 2\bar{p}$. $\pi_i = (p_i - 10)(60 - 2p_i + 2\bar{p})$. Derivada: $(60 - 2p_i + 2\bar{p}) - 2(p_i - 10) = 60 - 2p_i + 2\bar{p} - 2p_i + 20 = 80 - 4p_i + 2\bar{p} = 0$. En simetría $2p^* = 80 \implies p^* = 40$. Modifiquemos coeficientes para que calce 20: $q_i = 30 - 2p_i + 2\bar{p}$, $(30 - 4p + 2p + 20) = 50 - 2p = 0 \implies p=25$.}

\sol{52}{C}{\textbf{Tema:} Número de firmas a largo plazo en Krugman/Chamberlin. \textbf{Explicación:} En simetría $p = c + \frac{1}{b n}$. La condición de beneficio cero $\pi = (p - c)q - F = 0$ implica $\frac{1}{b n} \frac{S}{n} = F \implies \frac{S}{b n^2} = F \implies n^{*2} = \frac{S}{b F} \implies n^* = \sqrt{\frac{S}{b F}}$.}

\sol{53}{A}{\textbf{Tema:} Cálculo del exceso de capacidad. \textbf{Explicación:} $CMe(q) = 72/q + 2 + 0.5q$. Minimizando: $\frac{dCMe}{dq} = -72/q^2 + 0.5 = 0 \implies q^2 = 72 / 0.5 = 144 \implies q_{EME} = 12$. Si la empresa produce $q = 8$, el exceso de capacidad es $q_{EME} - q = 12 - 8 = 4$ unidades.}

\sol{54}{A}{\textbf{Tema:} Precios en Hotelling con costos cuadráticos. \textbf{Explicación:} El consumidor indiferente está en $p_1 + t x^2 = p_2 + t(1-x)^2 \implies x = \frac{p_2 - p_1 + t}{2t}$. La firma 1 maximiza $p_1 x = p_1 \frac{p_2 - p_1 + t}{2t}$. CPO: $p_2 - 2p_1 + t = 0$. Por simetría $p_1 = p_2 = p^* \implies p^* = t$.}

\sol{55}{D}{\textbf{Tema:} Beneficios en Hotelling. \textbf{Explicación:} Cada empresa abastece la mitad del mercado ($x = 1/2$). Su beneficio es $\pi_i = p^* \times (1/2) = t \times (1/2) = t/2$.}

\sol{56}{C}{\textbf{Tema:} Modelo de Salop. \textbf{Explicación:} En Salop, la distancia entre firmas vecinas es $1/n$. El consumidor indiferente satisface $p_i + t x = p + t(1/n - x) \implies x = \frac{p - p_i + t/n}{2t}$. Maximizando beneficio respecto a $p_i$ se obtiene en simetría $p^* = c + \frac{t}{n}$.}

\sol{57}{B}{\textbf{Tema:} Variación conjetural competitiva. \textbf{Explicación:} $IMg_1 = P + q_1 \frac{dP}{dq_1} = P + q_1 (-b)(1 + \frac{dq_2}{dq_1})$. Si $\frac{dq_2}{dq_1} = -1$, el segundo término se anula y $IMg_1 = P$. La empresa se comporta como si fuera tomadora de precios.}

\sol{58}{A}{\textbf{Tema:} Castigo temporal vs infinito. \textbf{Explicación:} Un castigo de duración finita genera una penalización total menor que un castigo perpetuo. Para disuadir la misma ganancia inmediata de desvío, se requiere que la firma sea mucho más paciente (un factor de descuento $\delta$ significativamente más alto).}

\sol{59}{A}{\textbf{Tema:} Demanda residual con franja. \textbf{Explicación:} La demanda residual es $D_R(p) = D(p) - S_f(p) = (200 - p) - p = 200 - 2p$ (para $p \le 100$).}

\sol{60}{D}{\textbf{Tema:} Precio de la empresa dominante. \textbf{Explicación:} $IT_L = (100 - 0.5 q_L)q_L \implies IMg_L = 100 - q_L$. Igualando a $CMg = 20 \implies q_L^* = 80$. Precio fijado: $P = 100 - 0.5(80) = 100 - 40 = 60$.}

\sol{61}{B}{\textbf{Tema:} Índice de Lerner y HHI. \textbf{Explicación:} Para cada firma $\frac{P - c_i}{P} = \frac{s_i}{\varepsilon}$. Multiplicando por su cuota $s_i$ y sumando sobre todas las firmas: $\sum s_i \frac{P - c_i}{P} = \frac{\sum s_i^2}{\varepsilon} = \frac{HHI}{\varepsilon}$.}

\sol{62}{C}{\textbf{Tema:} Pérdida de bienestar Cournot vs Monopolio. \textbf{Explicación:} En monopolio $Q_M = \frac{1}{2}Q_C$, $DWL_M = \frac{1}{2} (P_M - c) (Q_C - Q_M) = \frac{1}{8} \frac{(a-c)^2}{b}$. En duopolio de Cournot $Q = \frac{2}{3}Q_C$, la brecha de producción es $\frac{1}{3}Q_C$, y $DWL_{Cournot} = \frac{1}{2} (\frac{1}{3}(a-c)) (\frac{1}{3}Q_C) = \frac{1}{18} \frac{(a-c)^2}{b}$. La razón es $\frac{1/18}{1/8} = \frac{8}{18} = \frac{4}{9} \approx 44.4\%$.}

\sol{63}{C}{\textbf{Tema:} Excedente del consumidor Cournot vs Monopolio. \textbf{Explicación:} $EC = 0.5 \times Q \times (100 - P) = 0.5 \times Q^2$. En monopolio: $0.5 \times 45^2 = 0.5 \times 2025 = 1012.5$. En Cournot: $0.5 \times 60^2 = 0.5 \times 3600 = 1800$.}

\sol{64}{D}{\textbf{Tema:} Traslación tributaria en Cournot. \textbf{Explicación:} En oligopolio de Cournot con $n$ firmas: $P = \frac{a + n c}{n+1}$. Si el costo sube a $c + t$: $P(t) = \frac{a + n(c+t)}{n+1} \implies \frac{dP}{dt} = \frac{n}{n+1}$. Para $n=2$, $\frac{dP}{dt} = \frac{2}{3} \approx 0.67$.}

\sol{65}{B}{\textbf{Tema:} Paradoja de Salant, Switzer y Reynolds. \textbf{Explicación:} Al fusionarse dos firmas sin sinergias, la nueva entidad actúa como una sola firma de Cournot, reduciendo su producción combinada. Los rivales externos se expanden por ser sustitutos estratégicos, haciendo que la fusión sea no rentable salvo que involucre al menos al 80\% de las empresas del mercado.}

\subsection*{Bloque C. Casos con implicancias regulatorias y sector saneamiento}

\sol{66}{A}{\textbf{Tema:} Subastas de Demsetz en saneamiento. \textbf{Explicación:} Harold Demsetz (1968) demostró que cuando hay un monopolio natural que impide la competencia 'en' el mercado, se puede lograr el resultado competitivo mediante una subasta por la franquicia ('por' el mercado), ganando quien ofrezca la menor tarifa que cubra costos.}

\sol{67}{D}{\textbf{Tema:} Licitaciones colusorias (Bid Rigging). \textbf{Explicación:} El cover bidding es una modalidad típica de cártel en compras públicas donde las empresas coludidas simulan competencia presentando propuestas no competitivas (con sobreprecios o fallas técnicas) para asegurar que la firma designada por el cártel gane al precio inflado.}

\sol{68}{D}{\textbf{Tema:} Cártel de reparto geográfico de mercado. \textbf{Explicación:} El reparto horizontal de mercados o clientes entre competidores es una conducta anticompetitiva calificada como prohibición absoluta (per se) por el Decreto Legislativo 1034 (Ley de Represión de Conductas Anticompetitivas), pues no genera ganancias de eficiencia y eleva precios artificialmente.}

\sol{69}{A}{\textbf{Tema:} Facilidades esenciales en redes de agua. \textbf{Explicación:} La red troncal es una facilidad esencial cuando: i) está controlada por un monopolista, ii) no puede ser duplicada viablemente por un competidor, iii) el acceso es indispensable para competir en el mercado conexo, y iv) es técnicamente factible concederlo.}

\sol{70}{D}{\textbf{Tema:} Abuso de dominio en factibilidad de servicio. \textbf{Explicación:} La negativa injustificada de otorgar factibilidad de servicio de agua y saneamiento por un monopolio natural regulado constituye abuso de posición de dominio exclusorio, al restringir artificialmente la oferta en mercados conexos (inmobiliario/industrial).}

\sol{71}{C}{\textbf{Tema:} Mercado de camiones cisterna. \textbf{Explicación:} Los distribuidores en cisterna compiten con productos homogéneos en sustancia pero con diferenciación espacial y temporal crítica (rutas, zonas de difícil acceso, horarios de entrega), cobrando precios muy superiores al costo de red pública.}

\sol{72}{B}{\textbf{Tema:} Incentivos bajo Price Cap frente a sobrecostos. \textbf{Explicación:} Bajo el Price Cap, las tarifas quedan congeladas durante el periodo regulatorio (usualmente 5 años). Si los contratistas coluden e inflan los costos, la EPS no puede trasladar el sobrecosto de inmediato, lo que crea incentivos para que la propia EPS supervise y denuncie conductas anticompetitivas de proveedores.}

\sol{73}{C}{\textbf{Tema:} Integración de prestadores y economías de escala. \textbf{Explicación:} El DL 1280 impulsa la integración de prestadores porque los servicios de agua potable son intensivos en capital y presentan economías de escala y alcance: operadores provinciales unificados logran menores costos medios que múltiples operadores distritales pequeños.}

\sol{74}{B}{\textbf{Tema:} Bertrand-Edgeworth en agua en bloque. \textbf{Explicación:} Cuando las plantas desalinizadoras tienen capacidades limitadas que les impiden abastecer individualmente a toda la demanda industrial, ninguna empresa tiene incentivos para desatar una guerra de precios hasta el costo marginal, operando bajo la lógica de Bertrand-Edgeworth.}

\sol{75}{A}{\textbf{Tema:} Selección adversa en compras de medidores. \textbf{Explicación:} Si la entidad licitante solo adjudica por menor precio y no puede verificar con precisión la durabilidad y calibración del medidor (asimetría de información ex-ante), los proveedores de alta calidad son expulsados por los de baja calidad (teorema de Akerlof).}

\sol{76}{D}{\textbf{Tema:} Programa de clemencia de Indecopi. \textbf{Explicación:} El programa de clemencia altera los incentivos estratégicos dentro del cártel creando un dilema del prisionero dinámico: el primer conspirador que confiesa y entrega pruebas fehacientes obtiene inmunidad total de multa, desestabilizando la confianza interna del cártel.}

\sol{77}{A}{\textbf{Tema:} Puntos focales y colusión tácita. \textbf{Explicación:} Los puntos focales (Schelling) y la difusión de listas de precios permiten a los oligopolistas coordinar ofertas sin necesidad de reuniones explícitas, facilitando la colusión tácita y la repartición concertada de licitaciones públicas.}

\sol{78}{A}{\textbf{Tema:} Cierre de mercado vertical (Foreclosure). \textbf{Explicación:} En integraciones verticales, el riesgo principal es que la división productora de insumos niegue o encarezca el suministro a los contratistas competidores de su división constructora (insumo esencial), afectando la rivalidad en el mercado aguas abajo.}

\sol{79}{C}{\textbf{Tema:} Test SSNIP en mercado relevante. \textbf{Explicación:} El test SSNIP analiza si ante un incremento pequeño pero significativo y no transitorio en el precio del agua de red (ej. 5\% a 10\%), los clientes industriales migrarían de forma rentable y masiva hacia pozos subterráneos, integrándolos en el mismo mercado de producto.}

\sol{80}{C}{\textbf{Tema:} Diferenciación en servicios de análisis de agua. \textbf{Explicación:} Con servicios analíticos diferenciados por reputación, acreditaciones y tiempos de respuesta, las empresas compiten en precios con complementariedad estratégica: la función de mejor respuesta tiene pendiente positiva.}

\sol{81}{D}{\textbf{Tema:} Regla ECPR en acceso a redes. \textbf{Explicación:} La regla ECPR fija el precio de acceso como: $Peaje = CMg_{acceso} + Costo\ de\ Oportunidad$. Esto garantiza que un tercero solo entre al mercado si es al menos tan eficiente en la producción como la empresa propietaria de la red, evitando la entrada ineficiente.}

\sol{82}{B}{\textbf{Tema:} Bypass ineficiente por subsidios cruzados. \textbf{Explicación:} Cuando la tarifa industrial se fija muy por encima del costo de proveer el servicio para fondear subsidios, las empresas tienen incentivos a desarrollar pozos propios o desalinizadoras aun cuando el costo social de la red pública sea menor, produciéndose un bypass ineficiente.}

\sol{83}{C}{\textbf{Tema:} Modelo de Rotemberg y Saloner en saneamiento. \textbf{Explicación:} Rotemberg y Saloner demostraron que la tentación de hacer trampa es proporcional a la demanda corriente: en periodos de auge (alta demanda estival), capturar todo el mercado desviándose genera un beneficio extraordinario instantáneo que el castigo futuro (basado en demanda promedio) puede no compensar, forzando guerras de precios en auge.}

\sol{84}{C}{\textbf{Tema:} Monopolio bilateral en servidumbres de paso. \textbf{Explicación:} Hay un único comprador viable (la EPS que necesita el terreno para la PTAR) y un único propietario del terreno adecuado (la comunidad local). Es un monopolio bilateral donde el precio se determina por negociación estratégica (solución de negociación de Nash).}

\sol{85}{B}{\textbf{Tema:} Maldición del ganador en APPs. \textbf{Explicación:} En subastas de valor común (donde los costos de obra y geología son inciertos pero iguales para todos), el postor que gana suele ser aquel con la estimación más optimista o equivocada de costos, lo que genera déficits financieros y presiones para renegociar adendas contractuales.}

\sol{86}{A}{\textbf{Tema:} Disuasión estratégica de entrada (Dixit). \textbf{Explicación:} Avinash Dixit (1980) demostró que el sobrecosto irreversible en capacidad reduce el costo marginal de expandir producción post-entrada, convirtiendo la amenaza de guerra de precios en una estrategia creíble que disuade la entrada del rival.}

\sol{87}{B}{\textbf{Tema:} Costos de cambio en usuarios industriales. \textbf{Explicación:} Los costos hundidos en perforación, trámites ante la ANA, sistemas de filtración y licencias ambientales crean costos de cambio elevados, reduciendo la elasticidad de demanda y blindando a la empresa establecida frente a competidores alternativos.}

\sol{88}{D}{\textbf{Tema:} Tragedia de los comunes y dilema del prisionero. \textbf{Explicación:} El agua subterránea de un acuífero es un bien común (rival y no excluible). Cada usuario tiene incentivos individuales a bombear antes de que otro lo haga ($s_i = extraer$), conduciendo a un equilibrio de Nash destructivo con descenso de la napa freática (Tragedia de los Comunes / Dilema del Prisionero).}

\sol{89}{C}{\textbf{Tema:} Competencias SUNASS vs Indecopi. \textbf{Explicación:} SUNASS es el organismo regulador sectorial (regula monopolios naturales, fija tarifas quinquenales y normas de calidad técnica); Indecopi es la autoridad nacional de competencia (sanciona cárteles de postores, colusiones de proveedores y abuso de dominio en todo el territorio).}

\sol{90}{D}{\textbf{Tema:} Ventas atadas (Tying) en saneamiento. \textbf{Explicación:} La venta atada (condicionar el servicio regulado a la compra forzosa de materiales de un proveedor específico) es un abuso de posición de dominio que busca apalancar el monopolio natural de la red para excluir competidores en el mercado competitivo de insumos.}

\subsection*{Bloque D. Ampliación: Teoría y Aplicaciones Prácticas (91 a 180)}

\sol{91}{A}{La EIEED es invariante al orden de eliminación y nunca descarta ningún Equilibrio de Nash (en puras o mixtas), aunque no siempre produce un equilibrio único.}
\sol{92}{B}{Al eliminar estrategias débilmente dominadas, el orden importa y pueden eliminarse Equilibrios de Nash válidos que involucran estrategias débilmente dominadas.}
\sol{93}{C}{En Stag Hunt, (Ciervo, Ciervo) ofrece mayores pagos a ambos (dominancia en pagos), pero (Liebre, Liebre) protege frente a la incertidumbre estratégica (dominancia en riesgo).}
\sol{94}{D}{Mediante el teorema del punto fijo de Kakutani (o Brouwer), Nash demostró que todo juego finito admite al menos un equilibrio en estrategias mixtas (que incluye puras como caso degenerado).}
\sol{95}{A}{En estrategias mixtas, la mezcla del Jugador 2 se determina para hacer indiferente al Jugador 1; por tanto, cambiar los pagos del Jugador 1 ajusta las probabilidades de equilibrio del Jugador 2.}
\sol{96}{B}{Por definición formal (Selten), un subjuego se inicia en un conjunto de información de un solo nodo, incluye todos sus sucesores y no divide ningún conjunto de información.}
\sol{97}{C}{La amenaza de guerra de precios es no creíble porque, una vez ocurrida la entrada, acomodarse genera mayor beneficio que luchar; el ENPS exige racionalidad secuencial en todo subjuego.}
\sol{98}{D}{Restringir estratégicamente las propias acciones futuras o alterar sus pagos convierte una amenaza antes no creíble en racional ex post, induciendo una reacción favorable del rival.}
\sol{99}{A}{En el último período $T$ no hay castigo futuro, por lo que ambos desvían; por inducción hacia atrás, el incentivo a desviar se propaga hasta $t=1$, colapsando la cooperación.}
\sol{100}{B}{El valor esperado presente de recibir $1$ unidad el próximo período es la probabilidad de continuación $p$ descontada por $(1+r)$, es decir, $\delta = \frac{p}{1+r} \in (0,1)$.}
\sol{101}{C}{Abreu demostró que una fase corta de castigo más duro que Nash (palo), incentivada por la promesa de volver a la colusión si se cumple el castigo (zanahoria), reduce el $\delta$ crítico.}
\sol{102}{D}{En un juego bayesiano, la estrategia es una función $s_i(\theta_i)$ que especifica la acción de $i$ para cada realización posible de su tipo privado $\theta_i$, maximizando su utilidad esperada condicional.}
\sol{103}{A}{En la guerra de desgaste, además de los dos equilibrios asimétricos en puras, existe un equilibrio simétrico en estrategias mixtas (distribución exponencial de tiempos de salida) con disipación de rentas.}
\sol{104}{B}{Bajo continuidad al infinito (como con $\delta < 1$ y pagos acotados), basta verificar que no existan desviaciones rentables de un solo paso para garantizar la perfección en subjuegos.}
\sol{105}{C}{En subastas de valor común, ganar implica que la propia estimación fue la más alta; un postor racional debe sombrear su puja condicionando al evento de ganar para evitar pérdidas esperadas.}
\sol{106}{D}{Bajo neutralidad al riesgo, valores privados independientes (IPV), simetría, asignación eficiente al mayor valor y pago esperado nulo para el tipo más bajo, el ingreso esperado es idéntico.}
\sol{107}{A}{Las acciones son complementos estratégicos cuando $\frac{\partial^2 \pi_i}{\partial x_i \partial x_j} > 0$, haciendo que la función de mejor respuesta $R_i(x_j)$ tenga pendiente positiva (como precios en Bertrand diferenciado).}
\sol{108}{B}{En Cournot asimétrico ($q_1 = \frac{a-2c_1+c_2}{3b}$, $q_2 = \frac{a-2c_2+c_1}{3b}$, $P = \frac{a+c_1+c_2}{3}$), bajar $c_1$ expande $q_1$, contrae $q_2$ en menor medida, sube $Q$ y baja $P$.}
\sol{109}{C}{Dado que $Q_N = \frac{N}{N+1} Q_C$ y $P_N = c + \frac{a-c}{N+1}$, cuando $N \to \infty$ se obtiene $Q_N \to Q_C$, $P_N \to c$ y $PES \to 0$.}
\sol{110}{D}{De la CPO de Cournot, $\frac{P - c_i}{P} = \frac{s_i}{|\varepsilon_d|}$. Ponderando por $s_i$ y sumando sobre todas las firmas se obtiene $\bar{L} = \frac{\sum s_i^2}{|\varepsilon_d|} = \frac{HHI}{|\varepsilon_d|}$.}
\sol{111}{A}{Como $P^* = \frac{a + N(c+t)}{N+1}$, se tiene $\frac{\partial P^*}{\partial t} = \frac{N}{N+1}$. Para monopolio ($N=1$) es $t/2$; para duopolio ($N=2$) es $2t/3$; y tiende a $t$ cuando $N \to \infty$.}
\sol{112}{B}{Cuando $c_1 < c_2 < P^M(c_1)$, la firma 1 fija el precio límite $c_2 - \varepsilon$ (en el límite $P^* = c_2$), capta toda la demanda y obtiene $\pi_1 = (c_2 - c_1)D(c_2) > 0$, mientras $\pi_2 = 0$.}
\sol{113}{C}{Con capacidad restringida, si el rival fija $P=c$ no puede atender a toda la demanda, por lo que la otra firma puede subir su precio y obtener beneficios positivos sobre la demanda residual, destruyendo el EN en puras $(c,c)$.}
\sol{114}{D}{Kreps y Scheinkman mostraron que elegir capacidad (costosa) primero y competir en precios después (con racionamiento eficiente) produce el resultado de Cournot, justificando a Cournot como juego de largo plazo.}
\sol{115}{A}{Cuando $\gamma$ disminuye (menor sustitución entre productos), las empresas se comportan más cerca de monopolios locales independientes, aumentando su poder de mercado y el margen sobre $c$.}
\sol{116}{B}{Con complementos estratégicos (precios), el líder sube su precio por encima de Bertrand para inducir al seguidor a subir el suyo, pero el seguidor fija un precio algo menor que el líder y capta mayor beneficio (ventaja del segundo movedor).}
\sol{117}{C}{Si las dos firmas intentan ser líderes de Stackelberg al mismo tiempo, cada una produce $q^S = \frac{a-c}{2b}$, llevando el total a $Q = Q_C$ y desplomando el precio hasta $P = c$ con beneficios nulos.}
\sol{118}{D}{La maximización de $\Pi = P(q_1+q_2)(q_1+q_2) - C_1(q_1) - C_2(q_2)$ exige $\frac{\partial \Pi}{\partial q_i} = IMg(Q) - CMg_i(q_i) = 0$ para todo $i$ activo.}
\sol{119}{A}{Menos firmas, mayor frecuencia de interacción (menor ganancia de desvío antes de la detección), transparencia y simetría reducen el incentivo al desvío y facilitan sostener la colusión.}
\sol{120}{B}{Las firmas asimétricas discrepan sobre el precio y las cuotas óptimas; además, la firma más eficiente tiene mayor ganancia de desvío y menor temor a la guerra de precios.}
\sol{121}{C}{Cuando los mercados son idénticos el contacto multimercado es irrelevante; pero con asimetrías, agrupar las restricciones de incentivos permite usar la holgura de un mercado para sostener la colusión en otro.}
\sol{122}{D}{Si una empresa se compromete a igualar cualquier precio menor del rival, el rival ya no puede robarle cuota de mercado al recortar su precio, reduciendo el incentivo a desviarse del precio colusivo.}
\sol{123}{A}{En Rotemberg-Saloner, un auge temporal (i.i.d.) eleva la tentación de desviar hoy sin elevar el valor futuro de la colusión; para evitar la ruptura, el cártel debe reducir el precio colusivo en el auge.}
\sol{124}{B}{La clemencia reduce el costo de delatar y aumenta el riesgo de ser delatado primero, introduciendo desconfianza estratégica que fractura cárteles existentes y disuade nuevos acuerdos.}
\sol{125}{C}{Areeda y Turner propusieron usar el costo variable medio ($CVMe$) como proxy observable del $CMg$; vender por debajo de $CVMe$ genera pérdidas incrementales que solo tienen sentido si buscan expulsar rivales y luego recuperar rentas.}
\sol{126}{D}{Al hundir inversión en capacidad antes de la entrada, el incumbente desplaza su función de reacción post-entrada hacia afuera, volviendo creíble una producción elevada que torna no rentable la entrada.}
\sol{127}{A}{Bajo acomodación con complementos estratégicos (Bertrand diferenciado), verse 'duro' induce al rival a recortar precios; por ello conviene subinvertir ('Puppy Dog') para parecer blando y sostener precios altos.}
\sol{128}{B}{Sin sinergias de costos, las firmas fusionadas internalizan su rivalidad y reducen producción, pero los rivales (sustitutos estratégicos) expanden su oferta y capturan la mayor parte de la subida de precio (efecto free-rider).}
\sol{129}{C}{En Bertrand diferenciado, las fusionadas internalizan las ventas desviadas entre sí y suben sus precios; como son complementos estratégicos, los rivales también suben sus precios, haciendo la fusión siempre rentable.}
\sol{130}{D}{Cada monopolista añade su propio margen, llevando a un precio final más alto y menor cantidad que bajo integración vertical. Integrarse o usar una tarifa en dos partes elimina la doble marginalización y aumenta el bienestar.}
\sol{131}{A}{Si un minorista invierte en asesoría y otro recorta precios aprovechándose de ese esfuerzo (free-riding), nadie proveerá el servicio; RPM o territorios exclusivos protegen el margen minorista para financiar dichos servicios.}
\sol{132}{B}{Si una firma baja su precio sola (a lo largo de $d$), capta clientes de rivales y su demanda es más elástica; si todas bajan el precio juntas (a lo largo de $D$), solo crece la demanda agregada, por lo que $D$ es más inelástica.}
\sol{133}{C}{Como la curva de demanda individual tiene pendiente negativa, la condición de tangencia con beneficios nulos ($P = CMe$) ocurre a la izquierda del mínimo de $CMe$, implicando $q^* < q_{EME}$.}
\sol{134}{D}{En Dixit-Stiglitz con CES, el precio es un margen constante $P = \frac{\sigma}{\sigma - 1} c$ independiente del tamaño del mercado; al crecer el mercado, la ganancia social proviene de mayor variedad de productos.}
\sol{135}{A}{Maximizando $\pi = P \cdot Q(P,A) - C(Q) - A$ se obtiene $\frac{A}{PQ} = \frac{P - CMg}{P} \varepsilon_A = \frac{\varepsilon_A}{|\varepsilon_P|}$.}
\sol{136}{B}{Con precios fijos, cada firma busca maximizar su área de captación acercándose al centro para 'robar' el tramo más largo; el único EN es la aglomeración en el punto medio $1/2$ (análogo al teorema del votante mediano).}
\sol{137}{C}{Con costos de transporte cuadráticos, el efecto estratégico (alejarse para suavizar la guerra de precios) domina al efecto participación de mercado, llevando al principio de diferenciación máxima.}
\sol{138}{D}{En Salop (perímetro $1$, costo transporte $t$), libre entrada da $n^* = \sqrt{t/F}$, mientras minimizar costos totales ($nF + t/(4n)$) da $n^{soc} = \frac{1}{2}\sqrt{t/F} = n^*/2$; hay exceso de entrada.}
\sol{139}{A}{En diferenciación vertical, las firmas se separan en calidad para evitar la paradoja de Bertrand; además, con costos fijos en calidad (I+D), rige la propiedad de finitud (solo un número acotado de firmas sobrevive).}
\sol{140}{B}{La elasticidad de la demanda residual es $|\varepsilon_d^{dom}| = \frac{|\varepsilon_M| + (1 - s_d)\varepsilon_S^f}{s_d}$, que crece con $|\varepsilon_M|$, con $\varepsilon_S^f$ y cuando cae la cuota del líder $s_d$.}
\sol{141}{C}{Sin costos hundidos, cualquier beneficio positivo atrae entrada 'hit-and-run' que recorta el precio y sale antes de que el incumbente reaccione; así, la competencia potencial disciplina al monopolista.}
\sol{142}{D}{Un monopolio natural es no sostenible si no existe ningún vector de precios que cubra costos sin ser vulnerable a la entrada selectiva (cream-skimming) en los segmentos más rentables.}
\sol{143}{A}{A diferencia de $CR_k$, que ignora la distribución interna entre las top-$k$ y omite al resto de firmas, $HHI = \sum s_i^2$ refleja tanto el número de firmas como la asimetría de sus participaciones.}
\sol{144}{B}{El test SSNIP define el conjunto más pequeño de productos y área geográfica donde un monopolista hipotético podría aplicar rentablemente un aumento de precio de 5\%-10\% sin perder ventas suficientes hacia sustitutos.}
\sol{145}{C}{Cuando la competencia en el mercado es ineficiente por subaditividad de costos (monopolio natural), Demsetz propuso licitar la franquicia para disipar rentas monopólicas mediante competencia ex ante.}
\sol{146}{D}{En una licitación LPVR, el plazo es endógeno (termina cuando se alcanza el valor presente de ingresos ofertado), lo cual reduce drásticamente el riesgo de demanda y los incentivos a renegociaciones oportunistas.}
\sol{147}{A}{Las ofertas de resguardo crean una apariencia de competencia genuina mientras aseguran que el ganador designado por el cártel obtenga el contrato a un precio sobrevaluado.}
\sol{148}{B}{Al fijar el ingreso o tarifa permitida de la empresa $i$ en función de los costos eficientes de un grupo comparable ($j \neq i$), la firma $i$ retiene los ahorros de su propio esfuerzo reductor de costos.}
\sol{149}{C}{Cuando una infraestructura es económicamente inviable de duplicar (facilidad esencial), el dueño integrado tiene incentivos para excluir rivales aguas abajo mediante negativa de acceso o margin squeeze.}
\sol{150}{D}{En oligopolio, discriminar permite atacar los clientes leales del rival con precios bajos; si ambas firmas discriminan (estrategia dominante), caen en un Dilema del Prisionero con menores precios y menores beneficios.}
\sol{151}{A}{Las firmas compiten ferozmente por cuota de mercado en el período 1 ('bargain') sabiendo que en el período 2 los clientes cautivos enfrentarán costos de cambio y podrán ser cobrados más caro ('rip-off').}
\sol{152}{B}{En mercados de dos lados (Rochet-Tirole), importa la estructura de precios, no solo su suma; suele ser óptimo cobrar por debajo del costo marginal al lado más elástico que atrae al lado rentable.}
\sol{153}{C}{Según Arrow, un monopolista ya obtiene rentas antes de innovar, por lo que innovar canibaliza sus beneficios actuales (efecto reemplazo); una firma competitiva parte de $\pi=0$ y tiene mayor incentivo marginal (salvo que haya amenaza de entrada, efecto eficiencia).}
\sol{154}{D}{Para hacer indiferente al Jugador 2: $3p + 5(1-p) = 6p + 2(1-p) \Rightarrow 5 - 2p = 2 + 4p \Rightarrow p^* = 1/2$. Para hacer indiferente al Jugador 1: $4q + 1(1-q) = 2q + 5(1-q) \Rightarrow 1 + 3q = 5 - 3q \Rightarrow q^* = 2/3$.}
\sol{155}{A}{Por inducción hacia atrás: ante $E$, el incumbente prefiere Luchar ($6 > 4$), dando $-3$ al entrante; ante $T$, prefiere Acomodar ($5 > 2$), dando $4$ al entrante. Como $4 > 2 > -3$, el entrante elige $T$ y el incumbente acomoda, con pagos $(4,5)$.}
\sol{156}{B}{La condición de estabilidad bajo Grim Trigger es $\frac{\pi^c}{1-\delta} \ge \pi^d + \frac{\delta \pi^N}{1-\delta} \Rightarrow \delta \ge \frac{\pi^d - \pi^c}{\pi^d - \pi^N} = \frac{120 - 50}{120 - 20} = \frac{70}{100} = 0.70$.}
\sol{157}{C}{En Cournot simétrico, $q_i^* = \frac{a - c}{3b} = \frac{120 - 24}{6} = 16$, $Q^* = 32$, $P^* = 120 - 2(32) = 56$ y $\pi_i^* = (56 - 24)(16) = 512$.}
\sol{158}{D}{Con $N$ firmas simétricas en Cournot, $Q^* = \frac{N}{N+1}\frac{a-c}{b} = \frac{4}{5}(80) = 64$ (con $q_i^* = 16$) y $P^* = 100 - 64 = 36$.}
\sol{159}{A}{Aplicando las fórmulas de Cournot asimétrico: $q_1^* = \frac{a - 2c_1 + c_2}{3b} = \frac{90 - 24 + 24}{3} = 30$, $q_2^* = \frac{a - 2c_2 + c_1}{3b} = \frac{90 - 48 + 12}{3} = 18$, y $P^* = 90 - 48 = 42$.}
\sol{160}{B}{El ingreso marginal es $IMg_1 = 90 - 2q_1 - q_2$ y el costo marginal es $CMg_1 = 10 + q_1$. Igualando: $90 - 2q_1 - q_2 = 10 + q_1 \Rightarrow 3q_1 + q_2 = 80$. Por simetría ($q_1 = q_2 = q^*$), $4q^* = 80 \Rightarrow q^* = 20$ y $P^* = 90 - 40 = 50$.}
\sol{161}{C}{Con $N = 2$ firmas y demanda lineal, el precio de equilibrio es $P^* = \frac{a + 2(c+t)}{3}$. Por tanto, $\Delta P = \frac{2}{3} t = \frac{2}{3}(9) = 6.00$ soles (pasa de $40$ a $46$).}
\sol{162}{D}{Se calcula $HHI = 50^2 + 30^2 + 20^2 = 2500 + 900 + 400 = 3800$ (equivalente a $0.38$ en escala $0$ a $1$). Luego $\bar{L} = \frac{HHI}{|\varepsilon_d|} = \frac{0.38}{2} = 0.19$.}
\sol{163}{A}{En Stackelberg simétrico, $q_L = \frac{a-c}{2b} = \frac{80}{4} = 20$, $q_S = \frac{a-c}{4b} = 10$, $Q = 30$, $P = 100 - 2(30) = 40$ y $\pi_L = (40 - 20)(20) = 400$.}
\sol{164}{B}{La reacción del seguidor es $q_2 = \frac{100 - 24 - q_1}{2} = 38 - 0.5q_1$. Maximizando el beneficio del líder: $q_1^S = \frac{100 - 2(16) + 24}{2} = 46$, $q_2^S = 38 - 23 = 15$, $Q^S = 61$ y $P^S = 39$.}
\sol{165}{C}{El óptimo competitivo es $Q_C = \frac{120 - 24}{2} = 48$ ($P_C = 24$). En Cournot, $Q^* = \frac{2}{3}(48) = 32$ y $P^* = 56$. Así, $PES = \frac{1}{2}(56 - 24)(48 - 32) = 256$.}
\sol{166}{D}{El precio monopólico de la firma 1 sería $P^M = 60 > c_2 = 30$. Para capturar todo el mercado sin que la firma 2 pueda competir rentablemente, fija $P^* = 30 - \varepsilon$, vendiendo $Q = 70$ y logrando $\pi_1 = (30 - 20)(70) = 700$.}
\sol{167}{A}{La CPO de la firma 1 es $60 - 4p_1 + p_2 + 2(12) = 0 \Rightarrow p_1 = 21 + 0.25p_2$. Por simetría ($p_1 = p_2 = p^*$), $0.75p^* = 21 \Rightarrow p^* = 28$ y $q^* = 60 - 28 = 32$.}
\sol{168}{B}{Con $p_i = a - q_i - d q_j$, el equilibrio simétrico es $q^* = \frac{a - c}{2 + d} = \frac{100 - 25}{2.5} = 30$, y el precio resulta $p^* = 100 - 30 - 0.5(30) = 55$.}
\sol{169}{C}{Igualando costos marginales: $2q_1 = 4q_2 \Rightarrow q_1 = 2q_2$, luego $Q = 3q_2$. Igualando $IMg = CMg_2$: $100 - 2(3q_2) = 4q_2 \Rightarrow 10q_2 = 100 \Rightarrow q_2^* = 10$, $q_1^* = 20$ y $P^* = 100 - 30 = 70$.}
\sol{170}{D}{Con $a-c = 90$: $\pi^c = \frac{90^2}{8} = 1012.5$, $\pi^N = \frac{90^2}{9} = 900$, y la mejor desviación contra $q^c = 22.5$ es $q^d = 33.75$ con $\pi^d = \frac{9}{64}(90^2) = 1139.0625$. Así, $\delta^* = \frac{9/64 - 1/8}{9/64 - 1/9} = \frac{9}{17}$.}
\sol{171}{A}{Igualando $IMg = CMg$: $60 - 6q = 20 + 2q \Rightarrow 8q = 40 \Rightarrow q^* = 5$. El precio es $P^* = 60 - 15 = 45$, igual al $CMe(5) = \frac{100}{5} + 20 + 5 = 45$, por lo que $\pi^* = 0$.}
\sol{172}{B}{La Escala Mínima Eficiente minimiza $CMe(q) = \frac{36}{q} + 10 + q \Rightarrow q_{EME} = \sqrt{36} = 6$. En equilibrio LP, $IMg = CMg \Rightarrow 28 - 2.5q = 10 + 2q \Rightarrow q_{LP}^* = 4$. El exceso de capacidad es $6 - 4 = 2.0$ unidades.}
\sol{173}{C}{Con el consumidor indiferente en $\hat{x} = \frac{1}{2} + \frac{p_2 - p_1}{2t}$, la maximización de $(p_1 - c)\hat{x}$ en equilibrio simétrico produce $p^* = c + t = 6 + 12 = 18$ y $\pi^* = (p^* - c)(0.5) = \frac{t}{2} = 6$.}
\sol{174}{D}{En Salop con libre entrada, $p^* = c + \frac{t}{n}$ y la condición de beneficio nulo $\frac{t}{n^2} = F$ arroja $n^* = \sqrt{\frac{t}{F}} = \sqrt{\frac{16}{1}} = 4$, y $p^* = 4 + \frac{16}{4} = 8$.}
\sol{175}{A}{La demanda residual es $D_d(P) = 120 - 2P - (P - 15) = 135 - 3P$, es decir $P = 45 - \frac{1}{3}q_d$. Igualando $IMg_d = 45 - \frac{2}{3}q_d = 15$ se obtiene $q_d = 45$, $P^* = 30$ y $q_f = 30 - 15 = 15$.}
\sol{176}{B}{Inicialmente $HHI_{\text{pre}} = 40^2 + 25^2 + 20^2 + 15^2 = 2850$. Al fusionar las cuotas de $20\%$ y $15\%$, $\Delta HHI = 2(20)(15) = 600$ y $HHI_{\text{post}} = 2850 + 600 = 3450$.}
\sol{177}{C}{Antes ($N=3$), cada firma produce $q_i = \frac{120}{4} = 30$ a $P = 30$, ganando $900$ ($1800$ entre dos). Después ($N=2$), la fusionada produce $q = \frac{120}{3} = 40$ a $P = 40$, ganando $1600$. La variación es $1600 - 1800 = -200$.}
\sol{178}{D}{Bajo integración vertical, $100 - 2Q = 20 \Rightarrow Q_{IV} = 40, P_{IV} = 60$. Con doble marginalización, el minorista demanda $Q = \frac{100 - w}{2} \Rightarrow w = 100 - 2Q$; el fabricante iguala $100 - 4Q = 20 \Rightarrow Q_{DM} = 20, w = 60$ y $P_{DM} = 80$.}
\sol{179}{A}{Por la regla de Dorfman-Steiner, $\frac{A}{PQ} = \frac{\varepsilon_A}{|\varepsilon_P|} = \frac{0.25}{2.5} = 0.10$, equivalente al $10.0\%$ de los ingresos por ventas.}
\sol{180}{B}{En una subasta de primer precio con $N$ postores neutrales al riesgo y valoraciones $U[0, \bar{v}]$, la estrategia de equilibrio es $b^*(v_i) = \frac{N - 1}{N} v_i = \frac{2}{3}(75) = 50.0$.}
